% ============================================================
%  第一章 绪论
% ============================================================
\chapter{绪论}

\section{研究背景与意义}

近年来，随着高等教育规模的持续扩大，高校学生社团数量与成员人数呈现出快速增长的态势。社团作为高校第二课堂的重要组成部分，在培养学生综合素质、丰富校园文化生活方面发挥着不可替代的作用。然而，传统的社团管理模式普遍依赖纸质表单、电话沟通与人工协调，存在以下突出问题：

\begin{enumerate}
  \item \textbf{信息孤岛严重}：各社团数据分散存储，缺乏统一的信息平台，学校管理部门难以实时掌握社团运营状况；
  \item \textbf{流程繁琐低效}：社团注册、活动审批、招募报名等流程均需线下办理，耗时耗力，且容易出现信息丢失；
  \item \textbf{数据难以追溯}：活动签到、财务收支等记录缺乏电子化存档，审计与统计工作困难；
  \item \textbf{通知触达率低}：依赖微信群、公告栏等非正式渠道发布通知，信息传递不及时、不准确。
\end{enumerate}

信息化建设是解决上述问题的根本途径。构建一套功能完善、安全可靠的高校社团管理系统，不仅能够显著提升社团管理效率，降低行政成本，还能为学生提供更便捷的参与渠道，促进校园文化的繁荣发展。本系统的研究与实现具有重要的现实意义与应用价值。

\section{国内外研究现状}

\subsection{国内研究现状}

国内高校信息化建设起步于20世纪90年代，经过多年发展，教务管理、图书馆管理等核心业务系统已较为成熟。然而，针对学生社团管理的专项信息化系统建设相对滞后。目前，部分高校采用通用OA系统或自行开发的简单管理平台，但普遍存在功能单一、用户体验差、缺乏实时交互能力等不足\cite{wang2021campus}。近年来，随着Spring Boot、Vue等现代技术栈的普及，基于微服务或模块化架构的校园管理系统研究逐渐增多，但将AI智能服务与社团管理深度融合的研究仍较为少见。

\subsection{国外研究现状}

国外高校在学生事务管理信息化方面起步较早，以美国为代表的高校普遍采用商业化的学生信息系统（SIS），如Banner、PeopleSoft等，功能全面但定制化程度低、部署成本高\cite{chen2020university}。近年来，随着开源技术的发展，基于云原生架构的轻量级管理系统逐渐受到关注，但针对中国高校社团管理特点的本土化研究仍有较大空间。

\subsection{现有系统的不足}

综合国内外研究现状，现有系统主要存在以下不足：
\begin{enumerate}
  \item 缺乏完整的社团全生命周期管理（从申请、审批到解散）；
  \item 实时交互能力弱，通知推送依赖轮询而非WebSocket；
  \item AI智能服务集成度低，缺乏个性化推荐与智能问答功能；
  \item 安全机制不完善，缺乏细粒度的RBAC权限控制与审计日志。
\end{enumerate}

\section{研究内容与目标}

\subsection{系统功能目标}

本系统旨在实现以下核心功能：
\begin{enumerate}
  \item 用户注册、登录、权限管理及个人信息维护；
  \item 社团创建申请、审批、成员管理及生命周期管理；
  \item 活动发布、报名、签到码生成与核验、数据导出；
  \item 招募批次管理、动态表单配置与多阶段审核；
  \item 实时通知推送与公告管理；
  \item 财务收支记录、审批与余额管理；
  \item 场地与设备资源申请及冲突检测；
  \item AI智能社团推荐与知识库问答服务。
\end{enumerate}

\subsection{技术目标}

\begin{enumerate}
  \item \textbf{高安全性}：基于JWT双令牌机制与RBAC权限模型，防范XSS、CSRF、SQL注入等常见攻击；
  \item \textbf{高性能}：利用Redis缓存、数据库索引优化，保证核心接口响应时间在200ms以内；
  \item \textbf{高可用}：完善的异常处理与优雅降级机制，系统可用性达到99\%以上；
  \item \textbf{可扩展}：模块化架构设计，各业务模块低耦合，便于后续功能扩展与微服务化演进。
\end{enumerate}

\subsection{研究范围界定}

本系统面向高校校内使用场景，用户群体包括普通学生、社团成员、社团管理员与系统管理员。系统部署于校内服务器，通过浏览器访问，不涉及移动端原生应用开发。

\section{论文结构安排}

本文共分七章，各章内容安排如下：

第一章为绪论，介绍研究背景、国内外现状、研究目标与论文结构。

第二章介绍系统所采用的相关技术，包括后端技术栈、前端技术栈、AI服务及基础设施。

第三章进行系统需求分析，包括功能性需求、非功能性需求及用例分析。

第四章阐述系统设计，涵盖总体架构、数据库设计、安全架构、接口设计及AI服务设计。

第五章描述系统的具体实现，包括后端核心模块、前端页面及AI服务的代码实现。

第六章介绍系统测试，包括单元测试、接口测试、安全测试与性能测试。

第七章对全文进行总结，并展望未来的改进方向。
